
Three research questions were addressed:

\textbf{RQ1 (Existence):} Do domain-specific H_d(B,t) signals discriminate compressed from healthy benchmarks with statistically significant and practically large effect sizes?

\textbf{RQ2 (Granger-Predictive Mechanism):} Does submission count accumulation temporally precede score variance compression, or merely co-occur with it?

\textbf{RQ3 (Cross-sectional Diagnostic):} Do H_d signals serve as reliable indicators of currently-compressed benchmarks on real panel data?

\subsubsection{4.1 Datasets}

\begin{table}[htbp]
\centering
\begin{tabular}{llllll}
\toprule
Dataset & Source & Benchmarks & Panel Rows & Domains & Used for \\
\midrule
PWC Panel (real) & HuggingFace `pwc-archive/evaluation-tables` & 466 & 6,938 & CV, NLP & RQ2, RQ3 \\
Synthetic Panel & Generated (calibrated to real statistics) & 60 (20 per domain × 3) & N/A & CV, NLP, Tabular & RQ1 \\
\bottomrule
\end{tabular}
\end{table}

The PWC panel spans 2018–2025. The PWC REST API was unavailable during this research (shut down July 2025); the HuggingFace archive was used as a substitute. The synthetic panel for H-E1 was calibrated to match the median $\hat{\sigma}_{\text{measurement}} = 0.3323$ and the empirical submission count distribution from the real panel.

\subsubsection{4.2 Baselines}

\textbf{Spearman correlation (contemporaneous).} Cross-sectional Spearman ρ between cumulative submission count and score variance; tests linear correlation as an alternative to the lagged Granger structure.

\textbf{Granger null (reverse direction).} Granger causality from score variance to submission count; used as a directionality falsification check.

\textbf{Score variance + slope (naive).} Linear slope of score improvement and score variance across trailing 8 quarters, without domain-specific H_d design.

\subsubsection{4.3 Implementation}

All experiments were implemented in Python 3.10+ using \texttt{statsmodels} 0.14+ (Granger/ADF), \texttt{scipy.stats} 1.11+ (Mann-Whitney U, Spearman ρ), \texttt{numpy} 1.24+ (block-bootstrap Kendall τ), \texttt{lifelines} 0.27+ (Kaplan-Meier), and \texttt{datasets} 2.14+ (archive loading). All computations ran on CPU. Fixed seed = 42; fixed date range 2018-01-01 to 2025-12-31; fixed top-k = 10.

\subsubsection{4.4 Evaluation Metrics}

\begin{itemize}
\item \textbf{Mann-Whitney U p-value and Cohen's d (RQ1).} Gate: p < 0.05 AND |d| > 0.5 in at least 2/3 domains.
\item \textbf{Granger causality p-value at lag 2 (RQ2).} Gate: p < 0.05.
\item \textbf{AUC for binary classification (RQ3).} Gate: AUC > 0.8.
\item \textbf{Compression rate (descriptive).} Fraction of 466 qualifying benchmarks exhibiting at least 1 compression event.
\end{itemize}

\subsubsection{4.5 Hypothesis Execution Status}

\begin{table}[htbp]
\centering
\begin{tabular}{llllll}
\toprule
Hypothesis & Gate & Status & Reason &  &  \\
\midrule
H-E1 (signal discriminability) & MUST_WORK & PASS & p < 0.0001, \ & d\ & > 5 in all domains (synthetic data) \\
H-M1 (submission → compression, Granger) & MUST_WORK & PASS & Granger p = 1.854e-05 at lag = 2 &  &  \\
H-M2 (temporal ordering, H_d precedes collapse) & SHOULD_WORK & FAIL & Only 1 collapse event detected; gate requires ≥ 20 &  &  \\
H-M3 (Cox PH model, C-index ≥ 0.70) & MUST_WORK & NOT EXECUTED & Blocked by H-M2 collapse operationalization failure &  &  \\
H-M4 (Kaplan-Meier lead time ≥ 12 months) & SHOULD_WORK & NOT EXECUTED & Blocked by H-M3 dependency &  &  \\
\bottomrule
\end{tabular}
\end{table}

